% Reference list in the proposal-template format (numeric labels, in order of first citation).
% Generated from references.bib.
\begin{thebibliography}{63}
    \bibitem{sezer2020} Sezer, O. B., Gudelek, M. U., \& Ozbayoglu, A. M. (2020). Financial time series forecasting with deep learning: A systematic literature review: 2005--2019. Applied Soft Computing, 90, 106181. \url{https://doi.org/10.1016/j.asoc.2020.106181}
    \bibitem{fischer2018} Fischer, T., \& Krauss, C. (2018). Deep learning with long short-term memory networks for financial market predictions. European Journal of Operational Research, 270(2), 654--669. \url{https://doi.org/10.1016/j.ejor.2017.11.054}
    \bibitem{tetlock2007} Tetlock, P. C. (2007). Giving content to investor sentiment: the role of media in the stock market. The Journal of Finance, 62(3), 1139--1168. \url{https://doi.org/10.1111/j.1540-6261.2007.01232.x}
    \bibitem{kristoufek2013} Kristoufek, L. (2013). BitCoin meets Google Trends and Wikipedia: Quantifying the relationship between phenomena of the Internet era. Scientific Reports, 3, 3415. \url{https://doi.org/10.1038/srep03415}
    \bibitem{lo2004} Lo, A. W. (2004). The adaptive markets hypothesis: market efficiency from an evolutionary perspective. The Journal of Portfolio Management, 30(5), 15--29.
    \bibitem{cont2001} Cont, R. (2001). Empirical properties of asset returns: stylized facts and statistical issues. Quantitative Finance, 1(2), 223--236. \url{https://doi.org/10.1080/713665670}
    \bibitem{pineau2021} Pineau, J., et al. (2021). Improving reproducibility in machine learning research. Journal of Machine Learning Research, 22(164), 1--20.
    \bibitem{bouthillier2021} Bouthillier, X., et al. (2021). Accounting for variance in machine learning benchmarks. In Proceedings of Machine Learning and Systems, vol. 3.
    \bibitem{fama1970} Fama, E. F. (1970). Efficient Capital Markets: A Review of Theory and Empirical Work. The Journal of Finance, 25(2), 383--417. \url{https://doi.org/10.2307/2325486}
    \bibitem{urquhart2016} Urquhart, A. (2016). The inefficiency of Bitcoin. Economics Letters, 148, 80--82. \url{https://doi.org/10.1016/j.econlet.2016.09.019}
    \bibitem{nadarajah2017} Nadarajah, S., \& Chu, J. (2017). On the inefficiency of Bitcoin. Economics Letters, 150, 6--9.
    \bibitem{gu2020} Gu, S., Kelly, B., \& Xiu, D. (2020). Empirical asset pricing via machine learning. The Review of Financial Studies, 33(5), 2223--2273. \url{https://doi.org/10.1093/rfs/hhaa009}
    \bibitem{hochreiter1997} Hochreiter, S., \& Schmidhuber, J. (1997). Long Short-Term Memory. Neural Computation, 9(8), 1735--1780. \url{https://doi.org/10.1162/neco.1997.9.8.1735}
    \bibitem{gers2000} Gers, F. A., Schmidhuber, J., \& Cummins, F. (2000). Learning to forget: Continual prediction with LSTM. Neural Computation, 12(10), 2451--2471.
    \bibitem{mcnally2018} McNally, S., Roche, J., \& Caton, S. (2018). Predicting the price of Bitcoin using machine learning. In 26th Euromicro International Conference on Parallel, Distributed and Network-based Processing (pp. 339--343).
    \bibitem{jang2018} Jang, H., \& Lee, J. (2018). An empirical study on modeling and prediction of Bitcoin prices with Bayesian neural networks based on blockchain information. IEEE Access, 6, 5427--5437.
    \bibitem{huang2021} Huang, X., Zhang, W., Tang, X., Zhang, M., Surbiryala, J., Iosifidis, V., Liu, Z., \& Zhang, J. (2021). LSTM Based Sentiment Analysis for Cryptocurrency Prediction. arXiv preprint arXiv:2103.14804. \url{https://arxiv.org/abs/2103.14804}
    \bibitem{abraham2018} Abraham, J., Higdon, D., Nelson, J., \& Ibarra, J. (2018). Cryptocurrency price prediction using tweet volumes and sentiment analysis. SMU Data Science Review, 1(3). Article 1.
    \bibitem{girsang2023} Girsang, A. S., \& Stanley (2023). Cryptocurrency Price Prediction Based on Social Network Sentiment Analysis Using LSTM-GRU and FinBERT. IEEE Access. \url{https://doi.org/10.1109/ACCESS.2023.3324535}
    \bibitem{bollen2011} Bollen, J., Mao, H., \& Zeng, X. (2011). Twitter mood predicts the stock market. Journal of Computational Science, 2(1), 1--8. \url{https://doi.org/10.1016/j.jocs.2010.12.007}
    \bibitem{granger1969} Granger, C. W. J. (1969). Investigating causal relations by econometric models and cross-spectral methods. Econometrica, 37(3), 424--438. \url{https://doi.org/10.2307/1912791}
    \bibitem{loughran2011} Loughran, T., \& McDonald, B. (2011). When is a liability not a liability? Textual analysis, dictionaries, and 10-Ks. The Journal of Finance, 66(1), 35--65. \url{https://doi.org/10.1111/j.1540-6261.2010.01625.x}
    \bibitem{hutto2014} Hutto, C. J., \& Gilbert, E. (2014). VADER: A Parsimonious Rule-Based Model for Sentiment Analysis of Social Media Text. In Proceedings of the International AAAI Conference on Web and Social Media, vol. 8 (pp. 216--225). \url{https://doi.org/10.1609/icwsm.v8i1.14550}
    \bibitem{devlin2019} Devlin, J., Chang, M.-W., Lee, K., \& Toutanova, K. (2019). BERT: Pre-training of deep bidirectional transformers for language understanding. In Proceedings of the 2019 Conference of the North American Chapter of the Association for Computational Linguistics: Human Language Technologies (pp. 4171--4186).
    \bibitem{araci2019} Araci, D. (2019). FinBERT: Financial Sentiment Analysis with Pre-trained Language Models. arXiv preprint arXiv:1908.10063. \url{https://arxiv.org/abs/1908.10063}
    \bibitem{seabe2025} Seabe, P. L., Moutsinga, C. R. B., \& Pindza, E. (2025). Sentiment-driven cryptocurrency forecasting: analyzing LSTM, GRU, Bi-LSTM, and temporal attention model (TAM). Social Network Analysis and Mining, 15, 52. \url{https://doi.org/10.1007/s13278-025-01463-6}
    \bibitem{nezhad2026} Mohammadi Nezhad, M., Rouhani, S., Mohammadi, N., \& Shahedi, A. (2026). A unified GRU model for cryptocurrency price prediction and harsh price movement detection using enhanced sentiment analysis. Scientific Reports, 16, 16064. \url{https://doi.org/10.1038/s41598-026-46271-w}
    \bibitem{tiwari2025} Tiwari, D., Bhati, B. S., Nagpal, B., Alturki, N., \& Bayisenge, L. (2025). Attention-augmented hybrid CNN-LSTM model for social media sentiment analysis in cryptocurrency investment decision-making. Scientific Reports, 15, 33201. \url{https://doi.org/10.1038/s41598-025-18245-x}
    \bibitem{chen2016} Chen, T., \& Guestrin, C. (2016). XGBoost: A Scalable Tree Boosting System. In Proceedings of the 22nd ACM SIGKDD International Conference on Knowledge Discovery and Data Mining (pp. 785--794). \url{https://doi.org/10.1145/2939672.2939785}
    \bibitem{grinsztajn2022} Grinsztajn, L., Oyallon, E., \& Varoquaux, G. (2022). Why do tree-based models still outperform deep learning on typical tabular data? In Advances in Neural Information Processing Systems, Datasets and Benchmarks Track.
    \bibitem{brodersen2010} Brodersen, K. H., Ong, C. S., Stephan, K. E., \& Buhmann, J. M. (2010). The balanced accuracy and its posterior distribution. In 20th International Conference on Pattern Recognition (pp. 3121--3124).
    \bibitem{sokolova2009} Sokolova, M., \& Lapalme, G. (2009). A systematic analysis of performance measures for classification tasks. Information Processing and Management, 45(4), 427--437. \url{https://doi.org/10.1016/j.ipm.2009.03.002}
    \bibitem{tashman2000} Tashman, L. J. (2000). Out-of-sample tests of forecasting accuracy: an analysis and review. International Journal of Forecasting, 16(4), 437--450. \url{https://doi.org/10.1016/S0169-2070(00)00065-0}
    \bibitem{hyndman2021} Hyndman, R. J., \& Athanasopoulos, G. (2021). \emph{Forecasting: Principles and Practice} (3rd ed.). OTexts, Melbourne. \url{https://otexts.com/fpp3/}
    \bibitem{bergmeir2012} Bergmeir, C., \& Ben{\'i}tez, J. M. (2012). On the use of cross-validation for time series predictor evaluation. Information Sciences, 191, 192--213. \url{https://doi.org/10.1016/j.ins.2011.12.028}
    \bibitem{lopezdeprado2018} L{\'o}pez de Prado, M. (2018). \emph{Advances in Financial Machine Learning}. Wiley, Hoboken, NJ.
    \bibitem{bailey2014} Bailey, D. H., Borwein, J. M., L{\'o}pez de Prado, M., \& Zhu, Q. J. (2014). Pseudo-mathematics and financial charlatanism: The effects of backtest overfitting on out-of-sample performance. Notices of the American Mathematical Society, 61(5), 458--471.
    \bibitem{harvey2016} Harvey, C. R., Liu, Y., \& Zhu, H. (2016). \ldots and the cross-section of expected returns. The Review of Financial Studies, 29(1), 5--68. \url{https://doi.org/10.1093/rfs/hhv059}
    \bibitem{efron1993} Efron, B., \& Tibshirani, R. J. (1993). \emph{An Introduction to the Bootstrap}. Chapman and Hall, New York.
    \bibitem{dietterich1998} Dietterich, T. G. (1998). Approximate statistical tests for comparing supervised classification learning algorithms. Neural Computation, 10(7), 1895--1923. \url{https://doi.org/10.1162/089976698300017197}
    \bibitem{kunsch1989} K{\"u}nsch, H. R. (1989). The jackknife and the bootstrap for general stationary observations. The Annals of Statistics, 17(3), 1217--1241. \url{https://doi.org/10.1214/aos/1176347265}
    \bibitem{spearman1904} Spearman, C. (1904). The proof and measurement of association between two things. The American Journal of Psychology, 15(1), 72--101. \url{https://doi.org/10.2307/1412159}
    \bibitem{theiler1992} Theiler, J., Eubank, S., Longtin, A., Galdrikian, B., \& Farmer, J. D. (1992). Testing for nonlinearity in time series: the method of surrogate data. Physica D: Nonlinear Phenomena, 58(1--4), 77--94. \url{https://doi.org/10.1016/0167-2789(92)90102-S}
    \bibitem{binancepublic} Binance (2026). Binance Public Data: Spot Market Klines. Historical BTCUSDT 1-day OHLCV archives and checksum documentation; accessed 15 October 2026. \url{https://github.com/binance/binance-public-data}
    \bibitem{leetaru2013} Leetaru, K., \& Schrodt, P. A. (2013). GDELT: Global Data on Events, Location and Tone, 1979--2012. In ISA Annual Convention.
    \bibitem{gdeltdata} The GDELT Project (2026). GDELT Data: Querying, Analyzing and Downloading. Accessed 15 October 2026. \url{https://www.gdeltproject.org/data.html}
    \bibitem{gdelttone} The GDELT Project (2015). Global Knowledge Graph Data Format Codebook, Version 2.1. \url{https://data.gdeltproject.org/documentation/GDELT-Global_Knowledge_Graph_Codebook-V2.1.pdf}
    \bibitem{wilkinson2016} Wilkinson, M. D., et al. (2016). The FAIR Guiding Principles for scientific data management and stewardship. Scientific Data, 3, 160018. \url{https://doi.org/10.1038/sdata.2016.18}
    \bibitem{binanceterms} Binance (2026). Vision Dataset Terms and Conditions. Accessed 15 October 2026. \url{https://github.com/binance/binance-public-data/blob/master/TERMS_AND_CONDITIONS.md}
    \bibitem{gdeltterms} The GDELT Project (2026). GDELT Terms of Use and Citation. Accessed 15 October 2026. \url{https://www.gdeltproject.org/about.html}
    \bibitem{paszke2019} Paszke, A., et al. (2019). PyTorch: An imperative style, high-performance deep learning library. In Advances in Neural Information Processing Systems, vol. 32 (pp. 8024--8035).
    \bibitem{pedregosa2011} Pedregosa, F., et al. (2011). Scikit-learn: Machine learning in Python. Journal of Machine Learning Research, 12, 2825--2830.
    \bibitem{harris2020} Harris, C. R., et al. (2020). Array programming with NumPy. Nature, 585, 357--362. \url{https://doi.org/10.1038/s41586-020-2649-2}
    \bibitem{mckinney2010} McKinney, W. (2010). Data structures for statistical computing in Python. In Proceedings of the 9th Python in Science Conference.
    \bibitem{srivastava2014} Srivastava, N., Hinton, G., Krizhevsky, A., Sutskever, I., \& Salakhutdinov, R. (2014). Dropout: A simple way to prevent neural networks from overfitting. Journal of Machine Learning Research, 15(56), 1929--1958.
    \bibitem{kingma2015} Kingma, D. P., \& Ba, J. (2015). Adam: A method for stochastic optimization. In 3rd International Conference on Learning Representations. arXiv:1412.6980.
    \bibitem{pascanu2013} Pascanu, R., Mikolov, T., \& Bengio, Y. (2013). On the difficulty of training recurrent neural networks. In Proceedings of the 30th International Conference on Machine Learning, Proceedings of Machine Learning Research, vol. 28 (pp. 1310--1318).
    \bibitem{dunn1961} Dunn, O. J. (1961). Multiple comparisons among means. Journal of the American Statistical Association, 56(293), 52--64. \url{https://doi.org/10.1080/01621459.1961.10482090}
    \bibitem{fawcett2006} Fawcett, T. (2006). An introduction to ROC analysis. Pattern Recognition Letters, 27(8), 861--874. \url{https://doi.org/10.1016/j.patrec.2005.10.010}
    \bibitem{guo2017} Guo, C., Pleiss, G., Sun, Y., \& Weinberger, K. Q. (2017). On calibration of modern neural networks. In Proceedings of the 34th International Conference on Machine Learning, Proceedings of Machine Learning Research, vol. 70 (pp. 1321--1330).
    \bibitem{politis1994} Politis, D. N., \& Romano, J. P. (1994). The stationary bootstrap. Journal of the American Statistical Association, 89(428), 1303--1313. \url{https://doi.org/10.1080/01621459.1994.10476870}
    \bibitem{diebold1995} Diebold, F. X., \& Mariano, R. S. (1995). Comparing Predictive Accuracy. Journal of Business and Economic Statistics, 13(3), 253--263.
    \bibitem{brier1950} Brier, G. W. (1950). Verification of forecasts expressed in terms of probability. Monthly Weather Review, 78(1), 1--3.
\end{thebibliography}
